% SOURCE: Chinese PDF viewer pp.184 (middle)--186 (end of PDF).
\section{Plane-Wave Solutions of the Free Dirac Equation}
The Dirac equation for a free particle is
\[
(\ii\gamma^\mu\partial_\mu-m)\psi(x)=0.
\]
Assume that positive- and negative-energy plane waves take the forms
\[
\psi^+(x)=u(k)e^{-\ii k^\mu x_\mu},\qquad
\psi^-(x)=v(k)e^{+\ii k^\mu x_\mu},\qquad
k_0=E=\sqrt{m^2+\vect k^{\,2}}.
\]
Here $u(k)$ and $v(k)$ are Dirac spinors. Substitution yields
\[
(\gamma^\mu k_\mu-m)u(k)=0,\qquad
(\gamma^\mu k_\mu+m)v(k)=0.
\]

First consider the solutions in the particle rest frame, i.e. for $k_\mu=(m,\vect0)$:
\[
(\gamma_0-I)u(m,\vect0)=0,\qquad
(\gamma_0+I)v(m,\vect0)=0.
\]
Choose
\[
\gamma_0=\begin{pmatrix}I&0\\0&-I\end{pmatrix},\qquad
\vect\gamma=\begin{pmatrix}0&\vect\sigma\\-\vect\sigma&0\end{pmatrix}.
\]
$I$ is the identity matrix and $\vect\sigma$ the Pauli matrices. Thus each of $u$ and $v$ has two independent solutions:
\[
\begin{aligned}
u_\uparrow(m,\vect0)&=\begin{pmatrix}1\\0\\0\\0\end{pmatrix},&
u_\downarrow(m,\vect0)&=\begin{pmatrix}0\\1\\0\\0\end{pmatrix},\\[5pt]
v_\uparrow(m,\vect0)&=\begin{pmatrix}0\\0\\1\\0\end{pmatrix},&
v_\downarrow(m,\vect0)&=\begin{pmatrix}0\\0\\0\\1\end{pmatrix}.
\end{aligned}
\]
Let the Lorentz transformations from the particle rest frame to a moving frame be
\[
u_\alpha(k)=S_u u_\alpha(m,\vect0),\qquad
v_\alpha(k)=S_v v_\alpha(m,\vect0),\qquad
\alpha=\uparrow,\downarrow.
\]
Substituting the equations satisfied by $u(k)$ and $v(k)$ gives
\[
(\gamma^\mu k_\mu-m)S_u u_\alpha(m,\vect0)=0,\qquad
(\gamma^\mu k_\mu+m)S_v v_\alpha(m,\vect0)=0.
\]
Note that
\[
(\gamma^\mu k_\mu-m)(\gamma^\mu k_\mu+m)
=(\gamma^\mu k_\mu+m)(\gamma^\mu k_\mu-m)=k^2-m^2=0.
\]
Therefore one may take
\[
S_u=A(\gamma^\mu k_\mu+m),\qquad
S_v=B(\gamma^\mu k_\mu-m),
\]
where $A$ and $B$ are constants. Hence
\begin{align*}
u_\alpha(k)
&=A(\gamma^\mu k_\mu+m)u_\alpha(m,\vect0)
=A\begin{pmatrix}(E+m)\varphi_\alpha(m,\vect0)\\
\vect\sigma\cdot\vect k\,\varphi_\alpha(m,\vect0)\end{pmatrix},\\
v_\alpha(k)
&=B(\gamma^\mu k_\mu-m)v_\alpha(m,\vect0)
=-B\begin{pmatrix}\vect\sigma\cdot\vect k\,\varphi_\alpha(m,\vect0)\\
(E+m)\varphi_\alpha(m,\vect0)\end{pmatrix}.
\end{align*}
Here $\varphi_\alpha$ is a two-component spinor, with
\[
\varphi_\uparrow(m,\vect0)=\begin{pmatrix}1\\0\end{pmatrix},\qquad
\varphi_\downarrow(m,\vect0)=\begin{pmatrix}0\\1\end{pmatrix}.
\]
To determine $A$ and $B$ from the normalization conditions, define
\[
\overline u_\alpha=u_\alpha^\dagger\gamma^0,\qquad
\overline v_\alpha=v_\alpha^\dagger\gamma^0.
\]
Taking Hermitian conjugates,
\begin{align*}
\overline u_\alpha(k)
&=\overline u_\alpha(m,\vect0)(\gamma^\mu k_\mu+m)A
=A\left(\varphi_\alpha^{\mathsf T}(m,\vect0)(E+m),
-\varphi_\alpha^{\mathsf T}(m,\vect0)\vect\sigma\cdot\vect k\right),\\
\overline v_\alpha(k)
&=\overline v_\alpha(m,\vect0)(\gamma^\mu k_\mu-m)B
=-B\left(-\varphi_\alpha^{\mathsf T}(m,\vect0)\vect\sigma\cdot\vect k,
\varphi_\alpha^{\mathsf T}(m,\vect0)(E+m)\right).
\end{align*}
Define the orthonormality relations
\[
\begin{aligned}
\overline u_\alpha(k)u_\beta(k)&=\delta_{\alpha\beta},\\
\overline v_\alpha(k)v_\beta(k)&=-\delta_{\alpha\beta},\\
\overline u_\alpha(k)v_\beta(k)&=0,\\
\overline v_\alpha(k)u_\beta(k)&=0.
\end{aligned}
\]
The constants are then
\[
\boxed{A=-B=\frac{1}{\sqrt{2m(m+E)}}.}
\]